\section{Discussion}
% \textbf{Impact on the VLDB Community.}  
% This work bridges the gap between large language model (LLM) reasoning and graph-structured data systems, offering a unified, multi-agent \textit{Graph-CoT} framework co-designed with system-level serving optimizations. \todo{The proposed classification–reasoning–action decomposition introduces a modular abstraction that can be seamlessly integrated into abstract long-memory agent scenarios, such as knowledge graph question answering, information retrieval and task planning.} Furthermore, our priority-based KV cache eviction policy generalizes to other multi-agent and retrieval-augmented LLM serving systems, providing a practical path toward high-throughput, low-latency reasoning at scale. We believe this work opens new opportunities for the VLDB community to explore \textit{AI–data system co-design}, enabling future research on scalable, structure-aware reasoning frameworks that unify database principles with modern foundation models.


\stitle{Impact on AI–Data System Co-Design.}
This work bridges LLM reasoning and graph data management by introducing a multi-agent \textit{Graph-CoT} framework co-designed with system-level serving optimizations. The proposed classification–reasoning–action decomposition offers a reusable modular abstraction for accelerating long-memory agents\cite{zhang2025survey}, graph-enhanced LLMs\cite{liu2024knowledge}, and retrieval-augmented systems\cite{peng2024graph} across tasks such as knowledge-graph QA, structured information extraction, and task planning. Moreover, our vertex-centric KV-cache reuse and priority-based eviction policies generalize to multi-agent and structured RAG serving pipelines, providing a practical path toward high-throughput, low-latency inference at scale. We believe this work opens new directions for research at the intersection of data systems and LLMs, advancing scalable, structure-aware reasoning systems.



\stitle{Limitations and Future Work.}  
Graph-CoT opens several promising directions for future research. We highlight two main limitations of the current work and the corresponding opportunities for improvement in future research.  
First, the accuracy of \projectname{} remains bounded by the capability of the underlying backbone LLM. Future research may explore techniques such as knowledge distillation from larger models, lightweight fine-tuning (e.g., adapter or LoRA), tighter retrieval–reasoning calibration, and structure-aware prompting to reduce dependence on raw model capacity. 
Second, the graph retrieval process can introduce significant latency, particularly in the absence of GPU acceleration, and it often constitutes the dominant system bottleneck. Future work could investigate optimizing the retrieval pipeline and overall system architecture to reduce latency and resource consumption while preserving reasoning quality.



